\section{All-dimensional Q screening before generic standard discovery}
\label{sec:generic-euler}

The convexity theorem retained from report 81 applies in every matching
nullity, but the first Euler compiler screened only planar sections.
This continuation extends the screen to the generic branch. It addresses
an actual gap in that integration: the retained source bank contains 177
sectors of nullity four and eleven of nullity five that were outside the
planar gate. This is a supplied-sector improvement, with the same missing
universal coverage obligations as Section~\ref{sec:sector-windows}.

\subsection{Q hyperplanes suffice even without a full-dimensional cone}

Let $K$ be an injective rational matching basis with $d$ parameter coordinates
and $k$ allowed Q coordinates. Write $a_i$ for its coordinate forms, so that
$q_i=a_i\cdot z$ and $C=\{z:a_i\cdot z\ge0\}$.
The cone is pointed: a vector and its negative can both be feasible only
when all Q coordinates vanish, which by injectivity forces $z=0$.
Discard zero forms and deduplicate the hyperplanes using exact primitive
integer normals. Their number $h_Q$ is at most $k$.

\begin{proposition}[Complete generic Q-corner screening]
Every nonzero extreme ray of $C$ is generated by a one-dimensional kernel
of $d-1$ independent coordinate hyperplanes. Enumerating all such subsets,
retaining the feasible orientation and identifying proportional Q vectors
therefore enumerates every Q extreme direction. This remains true if the
linear span of $C$ has dimension smaller than $d$.
\end{proposition}
\begin{proof}
For a point $z$ in the relative interior of an extreme ray, let $I$ be its
zero Q coordinates. If the common kernel of their forms had dimension
at least two, choose a direction $u$ in it independent of $z$. The other
coordinates of $z$ are strictly positive, so sufficiently small positive
$\epsilon$ makes both $z+\epsilon u$ and $z-\epsilon u$ feasible. Their
average is $z$, contradicting extremality. The common kernel thus has
dimension one, and a basis of its normal space supplies $d-1$ active forms.
Conversely, any nonzero feasible direction in such a one-dimensional
intersection lies in a one-dimensional face of the pointed cone.
This argument includes constraints vanishing on the whole feasible cone;
no relative-interior promise in the ambient matching space is required.
\end{proof}

The new internal iterator solves those small rational systems and computes
the resulting Q direction. It selects the sign using the first nonzero Q
coordinate, rejects a direction containing any negative coordinate, and
emits only the lexicographically first independent subset of each ray's
active hyperplanes. Column pivots in the transpose of those active normals
select that subset by exact elimination. This gives polynomial deterministic
deduplication work per candidate without a hash-table assumption or storage
of previously emitted Q vectors. The parameter vector itself need not be
rescaled to an integral normal surface: canonical Euler characteristic is
positively homogeneous, so its sign is unchanged by positive rescaling.
The exact grouped objective from Section~\ref{sec:euler-aggregate} applies
in arbitrary parameter dimension. As before, the first corner uses the
direct functional; another corner triggers one grouping compilation.

A nonpositive value at every Q direction proves the Euler exclusion by
convexity. No partial scan does so. Source preparation checks the actual
compact manifold links, and callback interruption propagates before any
exclusion is returned. Positive Euler still requires ordinary standard
rays and independently verified disc components.

\subsection{Two integration points and a shared allowance}

Updated radius-one/two windows now use this iterator above nullity three,
before allocating their standard matching geometry. Lower-nullity chart
and scalar branches retain their previous behavior. A positive first Q
corner avoids grouping; any positive corner resumes the existing complete
standard enumeration. Source projections and positive proof formats remain
the same. The window stage continues to default off.

The public supplied-sector discovery API has a separate integration above
nullity three in automatic mode. Its existing Q-ray producer first lifts
and examines Q corners, checking any positive corner for an essential disc.
A found disc returns its ordinary independently replayed source witness.
A completed nonpositive scan returns the existing Q-phase exhaustion
certificate, whose stronger statement excludes positive Euler throughout
that sector, even though the caller requested standard discovery. The
certificate describes the phase that was actually completed. Explicit ray
methods and complete enumeration APIs retain their original behavior.

If the complete Q phase has positive Euler but finds no disc, automatic
discovery resumes its complete standard phase. The two phases share the
configured candidate-basis allowance: the standard phase receives the
initial allowance minus the Q candidates already begun. Thus a cap reached
in either phase returns \code{INCONCLUSIVE}, without an exhaustion proof.
The existing independent checker does not import the new producer or Euler
compiler. A positive Q phase can add work before the fallback, and a fixed
cap can consequently stop sooner than in the preceding standard-only
strategy. Neither a Q corner nor positive Euler alone is promoted to a
recognition verdict.

\subsection{What the asymptotic bound improves}

The generic internal screen begins at most
\[
                         \binom{h_Q}{d-1},\qquad h_Q\le k,
\]
small rational linear-algebra candidates, with sorting and canonical-active
subset selection also charged. Each has polynomial bit cost in
the supplied source, basis and coefficient encoding. At most this many
feasible directions need Euler evaluation; source compilation and grouping
also have polynomial preparation cost. The complete screen therefore costs
\[
                  \operatorname{poly}(T+k+B)(1+k)^{d-1}
\]
in bits, charging coefficient growth through exact rational elimination.
By contrast, the standard potential arrangement can have $O(k^2)$
hyperplanes and $\binom{O(k^2)}{d-1}$ candidate intersections. A completed
nonpositive screen avoids that arrangement and all its full coordinate
outputs and component queries. For a positive sector, both costs remain;
no uniformly smaller worst-case bound for complete discovery is asserted.

Logarithmic matching nullity again gives a conditional quasi-polynomial
local bound. This is not a general unknot algorithm: it supplies neither
logarithmic nullity nor a bounded complete family of admissible sectors for
arbitrary knot exteriors. The global recognition bound remains exponential.

\subsection{Independent source evidence and complete-query measurements}
All 1,422 maintained tests pass in 248.449 seconds, and the 23 focused
integration tests pass in 11.398 seconds. The tests cover a positive generic
corner, interrupted exclusion, a completed Q-only negative proof, a shared
cap after positive but nonessential Q rays, and repeat enumeration with one
statistics dictionary. An attempted standard arrangement is forbidden in
the Q-only exclusion test, so the smaller proof reflects an avoided search
rather than a shorter transcript of a search already performed.

The complete higher-nullity audit covers all 188 sectors selected by the
retained report 81 selection rule: 177 have nullity four and eleven have
nullity five. Every new Q stream equals the frozen previous Q stream in
coordinates and order. Direct and grouped Euler values agree with fresh
native source coordinate analysis on every Q ray, including scaling by
$2^{90}+1$. The screen's sign decision also agrees with the retained complete
standard-coordinate source oracle. Twelve sectors are excluded; each has
four Q corners and fifteen standard rays. The other 176 continue their
ordinary positive branch. Fresh Regina standard enumerations for all four
source triangulations represented by the exclusions reproduce exactly the
fifteen compatible rays in every case. All twelve returned Q exclusion
proofs pass independent replay.

The complete supplied-sector benchmark uses five randomized paired rounds
and four arms per case, including identical-baseline A/A controls. Its six
cases complete 120 measured calls and 24 warmups. Each call builds its source
kernel, completes automatic discovery, independently replays its returned
witness or exhaustion proof, and serializes the result. The previous arm
freezes release \code{e0a2aa2ee}; it does not substitute an earlier delivered
implementation. Both arms agree on every status, and all configured work
finishes. These are supplied-source queries, not full knot-recognition
measurements.

\begin{center}
\small
\begin{tabular}{lrrrrr}
\toprule
Source & previous ms & current ms & old/new & old A/A & new A/A\\
\midrule
\code{fibonacci\_lst\_09-d4} & 182.956 & 62.866 & 2.917 & 0.999 & 1.002\\
\code{fibonacci\_lst\_10-d4} & 228.834 & 83.590 & 2.745 & 1.001 & 0.986\\
\code{lst\_1\_8-d4} & 121.583 & 36.242 & 3.366 & 1.009 & 0.998\\
\code{solid\_torus\_sum\_s2xs1-d4} & 138.525 & 31.298 & 4.460 & 1.000 & 1.000\\
\code{cap\_1\_2\_3-d4} & 36.996 & 41.811 & 0.900 & 1.000 & 0.988\\
\code{cap\_3\_5\_3-d5} & 3502.927 & 3526.416 & 0.993 & 0.995 & 0.990\\
\bottomrule
\end{tabular}
\end{center}


The four nonpositive source cases have paired old/new ratios 2.917, 2.745,
3.366 and 4.460. Their candidate counts fall from 220 to 20 in the first
three cases and from 120 to four in the fourth. Each negative proof stores
four Q records instead of fifteen standard records. The A/A controls stay
within about two percent of parity. The positive nonessential nullity-four
control instead has ratio 0.900, about an eleven-percent slowdown, with
120 versus 124 candidates. The nullity-five control is near parity at
0.993, with 12,650 versus 12,655 candidates. Both positive controls finish
\code{NO\_VERTEX\_DISC\_IN\_SECTOR}; they demonstrate the cost of a Q phase
that must resume the complete standard search. No uniform speedup is claimed.


The 85-diagram native audit completes in 164.808 seconds. Every disabled
output, enabled verdict and callback work count agrees exactly with the
preceding release. All 31 positives remain; seven source-disc proofs receive
fresh Regina confirmation. The twenty-two caps and thirty-two completed
local misses also remain unchanged. None of these diagram windows enters
the new generic Q arrangement: they exercise the retained lower-nullity
branches. Thus the higher-nullity source-bank gains are not attributed to
improved coverage or work counts on this diagram corpus.

The diagram timing comparison completes 200 measured calls and forty
warmups over five native queries and five complete recognition calls,
including complete fallback when the native query is inconclusive. Fresh
input records, source and cache preparation, proof replay, configured work
and serialized output are charged. Every status and positive certificate
agrees across arms; none of these calls enters generic Q screening.

\textbf{Complete native queries with radius-two windows.}
\begin{center}
\small
\begin{tabular}{lrrrrr}
\toprule
Input & previous ms & current ms & old/new & old A/A & new A/A\\
\midrule
empty circle & 5.923 & 6.116 & 0.998 & 0.989 & 1.025\\
optimized positive & 57.609 & 59.896 & 0.988 & 0.990 & 0.976\\
genus one miss & 182.137 & 186.831 & 0.970 & 0.997 & 0.961\\
trefoil & 319.788 & 300.650 & 1.064 & 1.033 & 0.969\\
figure eight & 538.644 & 535.883 & 1.011 & 1.007 & 0.984\\
\bottomrule
\end{tabular}
\end{center}

\textbf{Complete recognition with radius-two windows enabled.}
\begin{center}
\small
\begin{tabular}{lrrrrr}
\toprule
Input & previous ms & current ms & old/new & old A/A & new A/A\\
\midrule
empty circle & 0.018 & 0.018 & 1.006 & 0.960 & 0.951\\
optimized positive & 51.192 & 54.268 & 0.937 & 1.000 & 1.071\\
genus one miss & 195.579 & 211.906 & 0.942 & 1.021 & 1.047\\
trefoil & 326.238 & 335.769 & 0.998 & 0.996 & 0.990\\
figure eight & 580.540 & 575.658 & 1.017 & 1.036 & 1.028\\
\bottomrule
\end{tabular}
\end{center}


These measurements do not establish a stable diagram-recognition gain.
In particular the genus-one recognition ratio is 0.942 while its new A/A
control is 1.047; the figure-eight controls are about 1.036 and 1.028.
Variation on these unchanged lower-nullity paths is not attributed to the
new generic algorithm, and the source-query gains are not generalized to
these diagrams. No reliable wall-clock gain is claimed for these calls.

Runtime release \code{71c8f8321}, the tests, source-sector audit, complete
source timings, 85-diagram audit and complete-recognition timings are
retained under \code{data/generic-euler-*}. From \code{fast/}, use
the module \code{normal\_orbit\_research.generic\_euler}, with
\code{sector-audit} and \code{--fresh-regina}, or its
\code{sector-benchmark}, \code{audit} and \code{benchmark} modes.
Publication review checks all 607 captured runtime pins, 605 preceding
release evidence pins, all timing outcomes, twelve source exclusions and
all 31 saved positives with the new Q iterator, Euler compiler and producers
disabled. Eight legacy generic proofs also receive fresh/context replay.
This verifies source exclusions and ordinary disc witnesses, without a new
global recognition theorem or a universal sector coverage claim.

